\documentclass[10pt,brazil,a4paper]{exam}
\usepackage[utf8]{inputenc}
%\usepackage[portuguese]{babel}
\usepackage[brazilian]{babel}
\usepackage[table,xcdraw]{xcolor}

%sudo apt install linuxbrew-wrapper
%sudo apt install texlive-extra-utils
%sudo apt install texlive-font-utils

\usepackage[pdftex]{graphicx}
\usepackage[dvips]{graphicx}
\usepackage{epstopdf}
\graphicspath{{./tmp/}}
\DeclareGraphicsExtensions{.pdf,.jpeg,.png,.eps}
\usepackage[outdir=./]{epstopdf}

\usepackage{hyperref} % include fig um http

\usepackage{listings} % para códigos

%\usepackage[top=.4in, bottom=.1in, left=1in, right=.4in]{geometry}
\textheight=6in
\usepackage{tikz}
\usepackage{enumitem}
%\usepackage[shortlabels]{enumerate}
\usepackage{multirow}
\usepackage{amsmath}
\usepackage{changepage,ifthen}
\usepackage{verbatim}
\usepackage{tabularx}
\usepackage{amsfonts}
\usepackage{multicol}

%\usepackage[T1]{fontenc}
%\usepackage{times}         
%\usepackage{multido}  % border
%\usepackage{pst-barcode}

\setlength{\textwidth}{185mm}
\setlength{\oddsidemargin}{-0.5in}
\setlength{\evensidemargin}{0in}
\setlength{\columnsep}{8mm}
\setlength{\topmargin}{-28mm}
\setlength{\textheight}{282mm}
\setlength{\itemsep}{0in}

\newcommand*\varhrulefill[1][0.4pt]{\leavevmode\leaders\hrule height#1\hfill\kern0pt}
\def\drawLines#1{{\color{lightgray}\foreach \x in {1,...,#1}{\par\vspace{2mm}\noindent\hrulefill}}}       
% cyan
\newcommand{\qid}[1]{\textcolor{white}{{\tiny #1}\hspace{0mm}}} %question identification (invisible text #1)
\usepackage[table,xcdraw]{xcolor}

#\renewcommand{\thefootnote}{\color{white}{\faGears}}
\renewcommand{\thefootnote}{\fnsymbol{footnote}}

\definecolor{bubbles}{rgb}{0.71, 1.0, 1.0}

\begin{document}
\pagestyle{empty}

\pagenumbering{gobble}
%\newgeometry{bottom=0.1cm}
\makeatletter\renewcommand*\cleardoublepage{\ifodd\c@page \else\hbox{}\newpage\fi}
\makeatother
\cleardoublepage
\newpage

% drawCircles
\vspace{-5mm}
\leavevmode\put(-11,0){\color{black}\circle*{15}}\hspace{-0mm}
\leavevmode\put(509,0){\color{black}\circle*{15}}
        \vspace{-1mm}\hspace{5mm}\begin{table}[h]
\begin{tabular}{|l|p{10.8cm}|l}
 \cline{1-2}\multirow{7}{*}{\vspace{8mm}\includegraphics[width=2cm]{./figs/webMCTest1}} 
&\textbf{Instituto Exemplo} 
              &\multirow{7}{*}{\hspace{-2mm}\includegraphics[scale=1.9000000000000001]{./tmp/QRCode_221_90_11201811685.eps}} \\ 
&\textbf{Curso Exemplo}                        & \\ 
&\textbf{Disciplina:} Disciplina Exemplo \hfill                        & \\ 
&\textbf{Prof.:} Francisco Zampirolli   & \\ 
&\textbf{Turma:} CE-teste
 \hfill \textbf{Sala:} 123
  & \\ 
&\textbf{Exame:} CE-exame teste \hfill \textbf{Data:} 20-10-2019           & \\ 
 \cline{1-2}\multicolumn{2}{|l|}{}      & \\ 
\multicolumn{2}{|l|}{\textbf{Ass.: }\rule{5cm}{0.1pt}} & \\ 
 \multicolumn{2}{|l|}{\textbf{\small{Estudante:}} ACOLON CAPONI DE CAIRES SILVA \hfill \textbf{\small{ID/RA:}} 11201811685}           & \\ 
 \cline{1-2}\end{tabular}
\end{table}
\vspace{-4.4mm}\hspace{-5mm}\footnote[2]{\color{lightgray}\textbf{webMCTest:} gerador e corretor de exames disponível para professores de instituições cadastradas em \textbf{\url{vision.ufabc.edu.br:8000}}}

\vspace{1.5mm}



 \hfill \tiny{{\color{red}2019-11-14 - 18:08:03\hspace{8mm}}}

\vspace{-6mm}\begin{center}\textbf{}\end{center}
\begin{center} 

 \vspace{-2mm} \ 
\begin{tikzpicture}[font=\tiny,fill=black]
  \foreach \numLab in {1,...,3} {
    \node[inner sep=3pt] at ({(\numLab-1+1) * 0.5},-0.1) {\numLab};
  }
  \foreach \letter/\line in {A/1,B/2,C/3,D/4,E/5} {
     \begin{scope}[xshift=0cm,yshift=-(\line)*5mm]
       \node at (0.1,0) {\letter};
       \foreach \letter/\position in {1,...,3} {
           \draw[fill=gray] ({\position * 0.5 - 0.2},-0.2) rectangle ({\position * 0.5+0.2},0.2); 
           \node[fill=white,draw,circle,inner sep=2.3pt] at ({\position * 0.5},0) {};
       }
     \end{scope}
  }
\end{tikzpicture}\hspace{1 mm}

 \ \ 
\end{center}

% drawCircles
\vspace{-5mm}
\leavevmode\put(-11,0){\color{black}\circle*{15}}\hspace{-0mm}
\leavevmode\put(509,0){\color{black}\circle*{15}}
        
%drawInstructions
\vspace{3mm}\noindent{\small\textbf{Instruções:}\vspace{-1mm}
\begin{enumerate}[label=(\alph*)]
\itemsep0pt\parskip0pt\parsep0pt
\item desligar o celular
\end{enumerate}}\normalsize

            \\\vspace{2mm}\hspace{-5mm}\noindent\textbf{Questões de Múltipla Escolha:}\vspace{2mm}


% QUESTOES DE MULTIPLA ESCOLHA



\hspace{-15mm}{\tiny {\color{white}\@1235}} \hspace{0mm} 1. A car moves on a road with an hourly function
$s=-4 + 6t$, where $s$ is
given in miles and $t$ in hours. The car passes
the mile 4 exactly at:

\vspace{0mm}

\begin{oneparchoices}
\choice \hspace{-2.0mm}{\tiny{\color{white}#1}}\hspace{2.0mm}2.33

\choice \hspace{-2.0mm}{\tiny{\color{white}\#0}}1.33

\choice \hspace{-2.0mm}{\tiny{\color{white}#4}}\hspace{2.0mm}-6.67

\choice \hspace{-2.0mm}{\tiny{\color{white}#2}}\hspace{2.0mm}-0.67

\choice \hspace{-2.0mm}{\tiny{\color{white}#3}}\hspace{2.0mm}0.33

\end{oneparchoices}\vspace{1mm}



\hspace{-15mm}{\tiny {\color{white}\@1236}} \hspace{0mm} 2. % Part I: description of a question 
Build a matrix of generic dimensions, for example, $62 \times 72$ whose elements $(i,j)$ are $((((i+1) * 19) + ((j+1) * 11)) \mod{100})$. Compute the sum of the entries of this matrix. Indexes $i$ of rows and $j$ of columns begin with $0$.

% Part II: we need to write cases here as a comment
\begin{comment}
{"input": [[75, 79, 19, 11], [74, 68, 19, 11], [62, 72, 19, 11]], "output": [[293250], [248944], [220900]]}
\end{comment}

% Part III: python code
\vspace{0mm}

\begin{oneparchoices}
\choice \hspace{-2.0mm}{\tiny{\color{white}\#0}}293250

\choice \hspace{-2.0mm}{\tiny{\color{white}#4}}\hspace{2.0mm}293248
\choice \hspace{-2.0mm}{\tiny{\color{white}#3}}\hspace{2.0mm}293249
\choice \hspace{-2.0mm}{\tiny{\color{white}#1}}\hspace{2.0mm}293240
\choice \hspace{-2.0mm}{\tiny{\color{white}#2}}\hspace{2.0mm}293253
\end{oneparchoices}\vspace{1mm}



\hspace{-15mm}{\tiny {\color{white}\@0070}} \hspace{0mm} 3. % texto da questão
O(a) \_\_\_\_\_\_ em Estrutura de Dados é também conhecido (a) como array Uni-dimensional. Assinale a alternativa que complete a lacuna.

\begin{verbatim}
 {"input": [["Uni", "lac-I"], ["entrada2"]], "output": [["Vetor", "lac-O"], ["saida2"]]}
\end{verbatim}

\begin{comment}
 {"input": [["Uni", "lac-I"], ["entrada2"]], "output": [["Vetor", "lac-O"], ["saida2"]]}
\end{comment}

\vspace{0mm}

\begin{oneparchoices}
\choice \hspace{-2.0mm}{\tiny{\color{white}\#0}}Vetor

\choice \hspace{-2.0mm}{\tiny{\color{white}#1}}\hspace{2.0mm}Matriz

\choice \hspace{-2.0mm}{\tiny{\color{white}#2}}\hspace{2.0mm}Pilha usando Ponteiro

\choice \hspace{-2.0mm}{\tiny{\color{white}#3}}\hspace{2.0mm}Árvore

\choice \hspace{-2.0mm}{\tiny{\color{white}#4}}\hspace{2.0mm}Grafo

\end{oneparchoices}\vspace{1mm}



% QUESTOES DE MULTIPLA ESCOLHA



% QUESTOES DE MULTIPLA ESCOLHA


 \ \ \ 
 \newpage
\makeatletter\renewcommand*\cleardoublepage{\ifodd\c@page \else\hbox{}\newpage\fi}
\makeatother
\cleardoublepage
\newpage

% drawCircles
\vspace{-5mm}
\leavevmode\put(-11,0){\color{black}\circle*{15}}\hspace{-0mm}
\leavevmode\put(509,0){\color{black}\circle*{15}}
        \vspace{-1mm}\hspace{5mm}\begin{table}[h]
\begin{tabular}{|l|p{10.8cm}|l}
 \cline{1-2}\multirow{7}{*}{\vspace{8mm}\includegraphics[width=2cm]{./figs/webMCTest1}} 
&\textbf{Instituto Exemplo} 
              &\multirow{7}{*}{\hspace{-2mm}\includegraphics[scale=1.9000000000000001]{./tmp/QRCode_221_90_11201721630.eps}} \\ 
&\textbf{Curso Exemplo}                        & \\ 
&\textbf{Disciplina:} Disciplina Exemplo \hfill                        & \\ 
&\textbf{Prof.:} Francisco Zampirolli   & \\ 
&\textbf{Turma:} CE-teste
 \hfill \textbf{Sala:} 123
  & \\ 
&\textbf{Exame:} CE-exame teste \hfill \textbf{Data:} 20-10-2019           & \\ 
 \cline{1-2}\multicolumn{2}{|l|}{}      & \\ 
\multicolumn{2}{|l|}{\textbf{Ass.: }\rule{5cm}{0.1pt}} & \\ 
 \multicolumn{2}{|l|}{\textbf{\small{Estudante:}} Acsa Santos Sousa \hfill \textbf{\small{ID/RA:}} 11201721630}           & \\ 
 \cline{1-2}\end{tabular}
\end{table}
\vspace{-4.4mm}\hspace{-5mm}\footnote[2]{\color{lightgray}\textbf{webMCTest:} gerador e corretor de exames disponível para professores de instituições cadastradas em \textbf{\url{vision.ufabc.edu.br:8000}}}

\vspace{1.5mm}



 \hfill \tiny{{\color{red}2019-11-14 - 18:08:03\hspace{8mm}}}

\vspace{-6mm}\begin{center}\textbf{}\end{center}
\begin{center} 

 \vspace{-2mm} \ 
\begin{tikzpicture}[font=\tiny,fill=black]
  \foreach \numLab in {1,...,3} {
    \node[inner sep=3pt] at ({(\numLab-1+1) * 0.5},-0.1) {\numLab};
  }
  \foreach \letter/\line in {A/1,B/2,C/3,D/4,E/5} {
     \begin{scope}[xshift=0cm,yshift=-(\line)*5mm]
       \node at (0.1,0) {\letter};
       \foreach \letter/\position in {1,...,3} {
           \draw[fill=gray] ({\position * 0.5 - 0.2},-0.2) rectangle ({\position * 0.5+0.2},0.2); 
           \node[fill=white,draw,circle,inner sep=2.3pt] at ({\position * 0.5},0) {};
       }
     \end{scope}
  }
\end{tikzpicture}\hspace{1 mm}

 \ \ 
\end{center}

% drawCircles
\vspace{-5mm}
\leavevmode\put(-11,0){\color{black}\circle*{15}}\hspace{-0mm}
\leavevmode\put(509,0){\color{black}\circle*{15}}
        
%drawInstructions
\vspace{3mm}\noindent{\small\textbf{Instruções:}\vspace{-1mm}
\begin{enumerate}[label=(\alph*)]
\itemsep0pt\parskip0pt\parsep0pt
\item desligar o celular
\end{enumerate}}\normalsize

            \\\vspace{2mm}\hspace{-5mm}\noindent\textbf{Questões de Múltipla Escolha:}\vspace{2mm}


% QUESTOES DE MULTIPLA ESCOLHA



\hspace{-15mm}{\tiny {\color{white}\@1235}} \hspace{0mm} 1. A car moves on a road with an hourly function
$s=-4 + 6t$, where $s$ is
given in miles and $t$ in hours. The car passes
the mile 4 exactly at:

\vspace{0mm}

\begin{oneparchoices}
\choice \hspace{-2.0mm}{\tiny{\color{white}#1}}\hspace{2.0mm}2.33

\choice \hspace{-2.0mm}{\tiny{\color{white}\#0}}1.33

\choice \hspace{-2.0mm}{\tiny{\color{white}#4}}\hspace{2.0mm}-6.67

\choice \hspace{-2.0mm}{\tiny{\color{white}#2}}\hspace{2.0mm}-0.67

\choice \hspace{-2.0mm}{\tiny{\color{white}#3}}\hspace{2.0mm}0.33

\end{oneparchoices}\vspace{1mm}



\hspace{-15mm}{\tiny {\color{white}\@1236}} \hspace{0mm} 2. % Part I: description of a question 
Build a matrix of generic dimensions, for example, $62 \times 72$ whose elements $(i,j)$ are $((((i+1) * 19) + ((j+1) * 11)) \mod{100})$. Compute the sum of the entries of this matrix. Indexes $i$ of rows and $j$ of columns begin with $0$.

% Part II: we need to write cases here as a comment
\begin{comment}
{"input": [[75, 79, 19, 11], [74, 68, 19, 11], [62, 72, 19, 11]], "output": [[293250], [248944], [220900]]}
\end{comment}

% Part III: python code
\vspace{0mm}

\begin{oneparchoices}
\choice \hspace{-2.0mm}{\tiny{\color{white}\#0}}293250

\choice \hspace{-2.0mm}{\tiny{\color{white}#4}}\hspace{2.0mm}293248
\choice \hspace{-2.0mm}{\tiny{\color{white}#3}}\hspace{2.0mm}293249
\choice \hspace{-2.0mm}{\tiny{\color{white}#1}}\hspace{2.0mm}293240
\choice \hspace{-2.0mm}{\tiny{\color{white}#2}}\hspace{2.0mm}293253
\end{oneparchoices}\vspace{1mm}



\hspace{-15mm}{\tiny {\color{white}\@0070}} \hspace{0mm} 3. % texto da questão
O(a) \_\_\_\_\_\_ em Estrutura de Dados é também conhecido (a) como array Uni-dimensional. Assinale a alternativa que complete a lacuna.

\begin{verbatim}
 {"input": [["Uni", "lac-I"], ["entrada2"]], "output": [["Vetor", "lac-O"], ["saida2"]]}
\end{verbatim}

\begin{comment}
 {"input": [["Uni", "lac-I"], ["entrada2"]], "output": [["Vetor", "lac-O"], ["saida2"]]}
\end{comment}

\vspace{0mm}

\begin{oneparchoices}
\choice \hspace{-2.0mm}{\tiny{\color{white}\#0}}Vetor

\choice \hspace{-2.0mm}{\tiny{\color{white}#1}}\hspace{2.0mm}Matriz

\choice \hspace{-2.0mm}{\tiny{\color{white}#2}}\hspace{2.0mm}Pilha usando Ponteiro

\choice \hspace{-2.0mm}{\tiny{\color{white}#3}}\hspace{2.0mm}Árvore

\choice \hspace{-2.0mm}{\tiny{\color{white}#4}}\hspace{2.0mm}Grafo

\end{oneparchoices}\vspace{1mm}



% QUESTOES DE MULTIPLA ESCOLHA



% QUESTOES DE MULTIPLA ESCOLHA


 \ \ \ 
 \newpage
\makeatletter\renewcommand*\cleardoublepage{\ifodd\c@page \else\hbox{}\newpage\fi}
\makeatother
\cleardoublepage
\newpage

% drawCircles
\vspace{-5mm}
\leavevmode\put(-11,0){\color{black}\circle*{15}}\hspace{-0mm}
\leavevmode\put(509,0){\color{black}\circle*{15}}
        \vspace{-1mm}\hspace{5mm}\begin{table}[h]
\begin{tabular}{|l|p{10.8cm}|l}
 \cline{1-2}\multirow{7}{*}{\vspace{8mm}\includegraphics[width=2cm]{./figs/webMCTest1}} 
&\textbf{Instituto Exemplo} 
              &\multirow{7}{*}{\hspace{-2mm}\includegraphics[scale=1.9000000000000001]{./tmp/QRCode_221_90_11001816.eps}} \\ 
&\textbf{Curso Exemplo}                        & \\ 
&\textbf{Disciplina:} Disciplina Exemplo \hfill                        & \\ 
&\textbf{Prof.:} Francisco Zampirolli   & \\ 
&\textbf{Turma:} CE-teste
 \hfill \textbf{Sala:} 123
  & \\ 
&\textbf{Exame:} CE-exame teste \hfill \textbf{Data:} 20-10-2019           & \\ 
 \cline{1-2}\multicolumn{2}{|l|}{}      & \\ 
\multicolumn{2}{|l|}{\textbf{Ass.: }\rule{5cm}{0.1pt}} & \\ 
 \multicolumn{2}{|l|}{\textbf{\small{Estudante:}} Adan Alves Siqueira \hfill \textbf{\small{ID/RA:}} 11001816}           & \\ 
 \cline{1-2}\end{tabular}
\end{table}
\vspace{-4.4mm}\hspace{-5mm}\footnote[2]{\color{lightgray}\textbf{webMCTest:} gerador e corretor de exames disponível para professores de instituições cadastradas em \textbf{\url{vision.ufabc.edu.br:8000}}}

\vspace{1.5mm}



 \hfill \tiny{{\color{red}2019-11-14 - 18:08:03\hspace{8mm}}}

\vspace{-6mm}\begin{center}\textbf{}\end{center}
\begin{center} 

 \vspace{-2mm} \ 
\begin{tikzpicture}[font=\tiny,fill=black]
  \foreach \numLab in {1,...,3} {
    \node[inner sep=3pt] at ({(\numLab-1+1) * 0.5},-0.1) {\numLab};
  }
  \foreach \letter/\line in {A/1,B/2,C/3,D/4,E/5} {
     \begin{scope}[xshift=0cm,yshift=-(\line)*5mm]
       \node at (0.1,0) {\letter};
       \foreach \letter/\position in {1,...,3} {
           \draw[fill=gray] ({\position * 0.5 - 0.2},-0.2) rectangle ({\position * 0.5+0.2},0.2); 
           \node[fill=white,draw,circle,inner sep=2.3pt] at ({\position * 0.5},0) {};
       }
     \end{scope}
  }
\end{tikzpicture}\hspace{1 mm}

 \ \ 
\end{center}

% drawCircles
\vspace{-5mm}
\leavevmode\put(-11,0){\color{black}\circle*{15}}\hspace{-0mm}
\leavevmode\put(509,0){\color{black}\circle*{15}}
        
%drawInstructions
\vspace{3mm}\noindent{\small\textbf{Instruções:}\vspace{-1mm}
\begin{enumerate}[label=(\alph*)]
\itemsep0pt\parskip0pt\parsep0pt
\item desligar o celular
\end{enumerate}}\normalsize

            \\\vspace{2mm}\hspace{-5mm}\noindent\textbf{Questões de Múltipla Escolha:}\vspace{2mm}


% QUESTOES DE MULTIPLA ESCOLHA



\hspace{-15mm}{\tiny {\color{white}\@0070}} \hspace{0mm} 1. % texto da questão
O(a) \_\_\_\_\_\_ em Estrutura de Dados é também conhecido (a) como array Uni-dimensional. Assinale a alternativa que complete a lacuna.

\begin{verbatim}
 {"input": [["Uni", "lac-I"], ["entrada2"]], "output": [["Vetor", "lac-O"], ["saida2"]]}
\end{verbatim}

\begin{comment}
 {"input": [["Uni", "lac-I"], ["entrada2"]], "output": [["Vetor", "lac-O"], ["saida2"]]}
\end{comment}

\vspace{0mm}

\begin{oneparchoices}
\choice \hspace{-2.0mm}{\tiny{\color{white}#2}}\hspace{2.0mm}Pilha usando Ponteiro

\choice \hspace{-2.0mm}{\tiny{\color{white}#4}}\hspace{2.0mm}Grafo

\choice \hspace{-2.0mm}{\tiny{\color{white}#1}}\hspace{2.0mm}Matriz

\choice \hspace{-2.0mm}{\tiny{\color{white}#3}}\hspace{2.0mm}Árvore

\choice \hspace{-2.0mm}{\tiny{\color{white}\#0}}Vetor

\end{oneparchoices}\vspace{1mm}



\hspace{-15mm}{\tiny {\color{white}\@1236}} \hspace{0mm} 2. % Part I: description of a question 
Build a matrix of generic dimensions, for example, $71 \times 70$ whose elements $(i,j)$ are $((((i+1) * 7) + ((j+1) * 23)) \mod{100})$. Compute the sum of the entries of this matrix. Indexes $i$ of rows and $j$ of columns begin with $0$.

% Part II: we need to write cases here as a comment
\begin{comment}
{"input": [[69, 60, 7, 23], [76, 62, 7, 23], [71, 70, 7, 23]], "output": [[205010], [233028], [245845]]}
\end{comment}

% Part III: python code
\vspace{0mm}

\begin{oneparchoices}
\choice \hspace{-2.0mm}{\tiny{\color{white}#2}}\hspace{2.0mm}205004
\choice \hspace{-2.0mm}{\tiny{\color{white}#4}}\hspace{2.0mm}205011
\choice \hspace{-2.0mm}{\tiny{\color{white}#1}}\hspace{2.0mm}205014
\choice \hspace{-2.0mm}{\tiny{\color{white}#3}}\hspace{2.0mm}205000
\choice \hspace{-2.0mm}{\tiny{\color{white}\#0}}205010

\end{oneparchoices}\vspace{1mm}



\hspace{-15mm}{\tiny {\color{white}\@1235}} \hspace{0mm} 3. A car moves on a road with an hourly function
$s=-3 + 7t$, where $s$ is
given in miles and $t$ in hours. The car passes
the mile 3 exactly at:

\vspace{0mm}

\begin{oneparchoices}
\choice \hspace{-2.0mm}{\tiny{\color{white}#4}}\hspace{2.0mm}-7.14

\choice \hspace{-2.0mm}{\tiny{\color{white}#1}}\hspace{2.0mm}1.86

\choice \hspace{-2.0mm}{\tiny{\color{white}\#0}}0.86

\choice \hspace{-2.0mm}{\tiny{\color{white}#2}}\hspace{2.0mm}-1.14

\choice \hspace{-2.0mm}{\tiny{\color{white}#3}}\hspace{2.0mm}-0.14

\end{oneparchoices}\vspace{1mm}



% QUESTOES DE MULTIPLA ESCOLHA



% QUESTOES DE MULTIPLA ESCOLHA


 \ \ \ 
 \newpage
\end{document}